\chapter{检索增强生成(RAG)}

\section{大语言模型的局限性}

在前面的章节中，我们学习了大型语言模型（LLM）的强大能力。然而，LLM 也存在一些固有的局限性：

\begin{itemize}
    \item \textbf{知识截止}：模型的训练数据有截止日期，无法获取最新信息
    \item \textbf{幻觉问题}：模型可能会生成看似合理但实际上错误的信息
    \item \textbf{领域知识}：对于特定领域的私有数据，模型无法访问
    \item \textbf{上下文长度}：模型的上下文窗口有限，无法处理超长文本
\end{itemize}

打个比方，LLM 就像一个博学的人，但他的知识停留在几年前，而且无法查阅最新的资料。如果你问他今天发生了什么，他只能凭记忆回答，可能会出错。

\section{什么是 RAG}

检索增强生成（Retrieval-Augmented Generation，简称 RAG）是一种结合检索和生成的技术。它的核心思想是：在生成回答之前，先从外部知识库中检索相关信息，然后将检索到的信息作为上下文提供给 LLM，让模型基于这些真实的信息生成回答。

这就像考试时允许翻书。你不需要记住所有知识，只需要知道如何查找相关信息，然后基于这些信息回答问题。

RAG 的基本流程如下：
\begin{enumerate}
    \item 用户提出问题
    \item 系统从知识库中检索相关文档
    \item 将问题和检索到的文档一起提供给 LLM
    \item LLM 基于检索到的信息生成回答
\end{enumerate}

\section{文本嵌入}

RAG 的第一步是将文本转换为向量表示，这个过程称为文本嵌入（Text Embedding）。

想象你在图书馆整理书籍。你会把相似主题的书放在同一个书架上。文本嵌入就是给每本书（文本）一个坐标（向量），使得内容相似的文本在向量空间中距离较近。

数学上，嵌入模型 $f$ 将文本 $t$ 映射到一个 $d$ 维向量：
\begin{equation*}
    f: t \rightarrow \mathbb{R}^d
\end{equation*}

例如：
\begin{verbatim}
文本1："机器学习是人工智能的一个分支"
文本2："深度学习属于机器学习领域"
文本3："今天天气很好"

嵌入后：
向量1：[0.2, 0.8, -0.3, ...]
向量2：[0.3, 0.7, -0.2, ...]  // 与向量1相似
向量3：[-0.5, 0.1, 0.9, ...]   // 与前两个不相似
\end{verbatim}

常用的相似度度量方法是余弦相似度：
\begin{equation*}
    \text{similarity} = \cos(\theta) = \frac{\mathbf{A} \cdot \mathbf{B}}{\|\mathbf{A}\| \|\mathbf{B}\|}
\end{equation*}

余弦值越接近 1，表示两个向量越相似。

\section{向量数据库}

检索到的文档需要存储在某个地方，这就是向量数据库的作用。向量数据库专门用于存储和检索高维向量。

打个比方，传统数据库就像电话簿，你通过名字（精确匹配）查找电话号码。向量数据库更像是一个智能地图，你说"我想找个附近好吃的中餐馆"，它会返回地理位置相近的选项。

常见的向量数据库包括：
\begin{itemize}
    \item FAISS（Facebook AI Similarity Search）
    \item Pinecone
    \item Milvus
    \item ChromaDB
\end{itemize}

向量数据库的核心功能是近似最近邻搜索（Approximate Nearest Neighbor, ANN）。当查询向量到来时，数据库快速找到与之最相似的 $k$ 个向量。

\section{RAG 的工作流程}

让我们通过一个完整的例子来理解 RAG 的工作流程。

假设你有一个公司内部的技术文档库，员工可以提问技术问题。

\subsection{索引阶段}

首先，需要对文档进行预处理和索引：

\begin{verbatim}
文档1：Python 的 GIL（全局解释器锁）确保同一时刻只有
       一个线程执行 Python 字节码。
文档2：使用 multiprocessing 模块可以绕过 GIL 的限制，
       实现真正的并行计算。
文档3：asyncio 是 Python 的异步编程框架，适合 I/O 密集型
       任务。
\end{verbatim}

处理步骤：
\begin{enumerate}
    \item \textbf{分块}：将长文档切分成较小的文本块（如 500 字一块）
    \item \textbf{嵌入}：使用嵌入模型（如 OpenAI 的 text-embedding-ada-002）将每个文本块转换为向量
    \item \textbf{存储}：将向量和对应的文本存储到向量数据库
\end{enumerate}

\subsection{查询阶段}

当用户提问时：

\begin{verbatim}
用户问题："如何在 Python 中实现并行计算？"
\end{verbatim}

处理步骤：
\begin{enumerate}
    \item \textbf{嵌入问题}：将用户问题转换为向量
    \item \textbf{检索}：在向量数据库中搜索最相似的文本块（如 top-3）
    \item \textbf{构建提示}：将问题和检索到的文本块组合成提示
    \item \textbf{生成回答}：LLM 基于提供的上下文生成回答
\end{enumerate}

构建的提示可能如下：

\begin{verbatim}
基于以下信息回答问题：

[检索到的文档块1]
使用 multiprocessing 模块可以绕过 GIL 的限制，
实现真正的并行计算。

[检索到的文档块2]
asyncio 是 Python 的异步编程框架，适合 I/O 密集型任务。

问题：如何在 Python 中实现并行计算？

回答：
\end{verbatim}

LLM 会基于检索到的信息生成：

\begin{verbatim}
在 Python 中实现并行计算，可以使用 multiprocessing 模块。
这个模块可以绕过 GIL（全局解释器锁）的限制，实现真正的
并行计算。如果需要处理 I/O 密集型任务，可以考虑使用
asyncio 异步编程框架。
\end{verbatim}

\section{RAG 的优势}

RAG 技术具有以下优势：

\begin{itemize}
    \item \textbf{知识实时更新}：只需更新知识库，无需重新训练模型
    \item \textbf{减少幻觉}：模型基于真实文档回答，可以追溯信息来源
    \item \textbf{领域特定}：可以构建特定领域的知识库，提供专业回答
    \item \textbf{成本效益}：无需微调模型，只需维护知识库
\end{itemize}

\section{RAG 的挑战}

尽管 RAG 很强大，但也面临一些挑战：

\begin{itemize}
    \item \textbf{检索质量}：如果检索到的文档不相关，会影响生成质量
    \item \textbf{分块策略}：如何合理切分文档是个技术活
    \item \textbf{上下文长度}：检索太多文档会超出模型的上下文窗口
    \item \textbf{延迟}：检索和生成两个步骤增加了响应时间
\end{itemize}

针对这些问题，研究者提出了许多改进方法，如：
\begin{itemize}
    \item \textbf{混合检索}：结合关键词检索和向量检索
    \item \textbf{重排序}：对检索结果进行二次排序
    \item \textbf{查询改写}：优化用户查询以提高检索效果
\end{itemize}
